\chapter{2009年陕西卷（文）}

\section{单选题}

\begin{problem}
设不等式 \(x^{2} - x \leqslant 0\) 的解集为 \(M\)，函数 \(f(x) = \ln (1 - |x|)\) 的定义域为 \(N\)，则 \(M \cap N\) 为（\quad）

\choices
    {\([0,1)\)}
    {\((0,1)\)}
    {\([0,1]\)}
    {\((-1,0]\)}
\end{problem}

\begin{answer}
A
\end{answer}

\begin{solution}
由 \(x^{2}-x\leqslant 0\) 得 \(x(x-1)\leqslant 0\)，所以 \(M=[0,1]\). 函数 \(\ln(1-|x|)\) 有意义要求 \(1-|x|>0\)，即 \(-1<x<1\)，所以 \(N=(-1,1)\). 因而 \(M\cap N=[0,1)\)，故选 A.
\end{solution}

\begin{problem}
若 \(\tan \alpha = 2\)，则 \(\frac{2\sin\alpha - \cos\alpha}{\sin\alpha + 2\cos\alpha}\) 的值为（\quad）

\choices
    {\(0\)}
    {\( \frac{3}{4} \)}
    {\(1\)}
    {\(\frac{5}{4}\)}
\end{problem}

\begin{answer}
B
\end{answer}

\begin{solution}
由 \(\tan\alpha=2\) 可知 \(\cos\alpha\neq 0\). 分子分母同除以 \(\cos\alpha\)，原式等于 \(\dfrac{2\tan\alpha-1}{\tan\alpha+2}=\dfrac{4-1}{2+2}=\dfrac{3}{4}\)，故选 B.
\end{solution}

\begin{problem}
函数 \( f(x) = \sqrt{2x - 4} \)（\(x \geqslant 4\)）的反函数为（\quad）

\choices
    {\(f^{-1}(x) = \frac{1}{2} x^{2} + 4\)（\(x \geqslant 0\)）}
    {\(f^{-1}(x) = \frac{1}{2} x^2 + 4\)（\(x \geqslant 2\)）}
    {\(f^{-1}(x) = \frac{1}{2} x^2 + 2\)（\(x \geqslant 0\)）}
    {\(f^{-1}(x) = \frac{1}{2} x^2 + 2\)（\(x \geqslant 2\)）}
\end{problem}

\begin{answer}
D
\end{answer}

\begin{solution}
设 \(y=f(x)=\sqrt{2x-4}\). 题设把原函数的定义域限定为 \(x\geqslant 4\)，所以其值域为 \(y\geqslant 2\). 由 \(y^{2}=2x-4\) 得 \(x=\dfrac{y^{2}}{2}+2\). 交换 \(x\)，\(y\)，并把原函数的值域作为反函数的定义域，得到 \(f^{-1}(x)=\dfrac{1}{2}x^{2}+2\)（\(x\geqslant 2\)），故选 D.
\end{solution}

\begin{problem}
过原点且倾斜角为 \(60^{\circ}\) 的直线被圆 \(x^{2} + y^{2} - 4y = 0\) 所截得的弦长为（\quad）

\choices
    {\(\sqrt{3}\)}
    {\(2\)}
    {\( \sqrt{6} \)}
    {\(2\sqrt{3}\)}
\end{problem}

\begin{answer}
D
\end{answer}

\begin{solution}
圆的方程可化为 \(x^{2}+(y-2)^{2}=4\)，所以圆心为 \(O(0,2)\)，半径为 \(2\). 直线过原点且倾斜角为 \(60^{\circ}\)，其方程为 \(y=\sqrt{3}x\)，即 \(\sqrt{3}x-y=0\). 圆心到直线的距离为 \(d=\dfrac{2}{\sqrt{3+1}}=1\). 弦长为 \(2\sqrt{2^{2}-1^{2}}=2\sqrt{3}\)，故选 D.
\end{solution}

\begin{problem}
某单位共有老、中、青职工 \(430\) 人，其中青年职工 \(160\) 人，中年职工人数是老年职工人数的 \(2\) 倍. 为了解职工身体状况，现采用分层抽样方法进行调查，在抽取的样本中有青年职工 \(32\) 人，则该样本中的老年职工人数为（\quad）

\choices
    {\(9\)}
    {\(18\)}
    {\(27\)}
    {\(36\)}
\end{problem}

\begin{answer}
B
\end{answer}

\begin{solution}
设老年职工人数为 \(x\)，则中年职工人数为 \(2x\). 由总人数得 \(x+2x+160=430\)，解得 \(x=90\). 分层抽样中各层的抽样比例相同，青年职工的抽样比例为 \(\dfrac{32}{160}=\dfrac{1}{5}\). 因而老年职工的抽样人数为 \(90\times\dfrac{1}{5}=18\)，故选 B.
\end{solution}

\begin{problem}
若 \( (1-2x)^{2009}=a_{0}+a_{1}x+\cdots+a_{2009}x^{2009} \)（\(x\in\R\)），则 \( \frac{a_{1}}{2}+\frac{a_{2}}{2^{2}}+\cdots+\frac{a_{2009}}{2^{2009}} \)（\quad）

\choices
    {\(2\)}
    {\(0\)}
    {\(-1\)}
    {\(-2\)}
\end{problem}

\begin{answer}
C
\end{answer}

\begin{solution}
令 \(x=0\)，由展开式得 \(a_{0}=1\). 令 \(x=\dfrac{1}{2}\)，则 \(0=(1-2\cdot\dfrac{1}{2})^{2009}=1+\dfrac{a_{1}}{2}+\dfrac{a_{2}}{2^{2}}+\cdots+\dfrac{a_{2009}}{2^{2009}}\). 因此题中所求的和为 \(-1\)，故选 C.
\end{solution}

\begin{problem}
“\(m > n > 0\)”是“方程 \(mx^{2} + ny^{2} = 1\) 表示焦点在 \(y\) 轴上的椭圆”的（\quad）

\choices
    {充分而不必要条件}
    {必要而不充分条件}
    {充要条件}
    {既不充分也不必要条件}
\end{problem}

\begin{answer}
C
\end{answer}

\begin{solution}
因为 \(m>n>0\)，方程可写为 \(\dfrac{x^{2}}{1/m}+\dfrac{y^{2}}{1/n}=1\). 其半轴长的平方分别为 \(1/m\) 和 \(1/n\)，而 \(m>n>0\) 等价于 \(1/n>1/m\)，所以较长半轴在 \(y\) 轴上，焦点在 \(y\) 轴上，故该条件充分. 反过来，若该方程表示焦点在 \(y\) 轴上的椭圆，则必有 \(m\)，\(n>0\)，且 \(1/n>1/m\)，从而 \(m>n\). 因此该条件也是必要条件，故选 C.
\end{solution}

\begin{problem}
在 \(\triangle ABC\) 中，\(M\) 是 \(BC\) 的中点，\(AM = 1\)，点 \(P\) 在 \(AM\) 上且满足 \(\overrightarrow{AP} = 2\overrightarrow{PM}\)，则 \(\overrightarrow{PA} \cdot (\overrightarrow{PB} + \overrightarrow{PC})\) 等于（\quad）

\choices
    {\( \frac{4}{9} \)}
    {\( \frac{4}{3} \)}
    {\(-\frac{4}{3}\)}
    {\(-\frac{4}{9}\)}
\end{problem}

\begin{answer}
D
\end{answer}

\begin{solution}
取 \(A\)，\(B\)，\(C\)，\(M\)，\(P\) 的位置向量分别为 \(\bs{a}\)，\(\bs{b}\)，\(\bs{c}\)，\(\bs{m}\)，\(\bs{p}\). 因为 \(M\) 是 \(BC\) 的中点，故 \(\bs{m}=\dfrac{\bs{b}+\bs{c}}{2}\). 又由 \(\overrightarrow{AP}=2\overrightarrow{PM}\) 得 \(\bs{p}-\bs{a}=2(\bs{m}-\bs{p})\)，所以 \(\bs{p}=\dfrac{\bs{a}+2\bs{m}}{3}\). 于是 \(\overrightarrow{PA}=\bs{a}-\bs{p}=\dfrac{2}{3}(\bs{a}-\bs{m})\)，且 \(\overrightarrow{PB}+\overrightarrow{PC}=\bs{b}+\bs{c}-2\bs{p}=2\bs{m}-2\bs{p}=-\dfrac{2}{3}(\bs{a}-\bs{m})\). 因而
\[
\overrightarrow{PA}\cdot(\overrightarrow{PB}+\overrightarrow{PC})=-\dfrac{4}{9}|\bs{a}-\bs{m}|^{2}=-\dfrac{4}{9}AM^{2}=-\dfrac{4}{9}
\]
故选 D.
\end{solution}

\begin{problem}
从 \(1\)，\(2\)，\(3\)，\(4\)，\(5\)，\(6\)，\(7\) 这七个数字中任取两个奇数和两个偶数，组成没有重复数字的四位数，其中奇数的个数为（\quad）

\choices
    {\(432\)}
    {\(288\)}
    {\(216\)}
    {\(108\)}
\end{problem}

\begin{answer}
A
\end{answer}

\begin{solution}
七个数字中有奇数 \(1\)，\(3\)，\(5\)，\(7\) 共 \(4\) 个，偶数 \(2\)，\(4\)，\(6\) 共 \(3\) 个. 先选出两个奇数和两个偶数，有 \(\mathrm C_{4}^{2}\mathrm C_{3}^{2}\) 种；选出的四个数字互不相同，排成四位数有 \(4!\) 种. 因而总数为 \(\mathrm C_{4}^{2}\mathrm C_{3}^{2}4!=6\times3\times24=432\)，故选 A.
\end{solution}

\begin{problem}
定义在 \(\R\) 上的偶函数 \(f(x)\) 满足：对任意的 \(x_{1}\)，\(x_{2} \in [0, +\infty)\)（\(x_{1} \neq x_{2}\)），有 \(\frac{f(x_2) - f(x_1)}{x_2 - x_1} < 0\). 则（\quad）

\choices
    {\(f(3) < f(-2) < f(1)\)}
    {\(f(1) < f(-2) < f(3)\)}
    {\(f(-2) < f(1) < f(3)\)}
    {\(f(3) < f(1) < f(-2)\)}
\end{problem}

\begin{answer}
A
\end{answer}

\begin{solution}
题设差商对任意不同的 \(x_{1}\)，\(x_{2}\in[0,+\infty)\) 都小于 \(0\)，所以 \(f\) 在 \([0,+\infty)\) 上严格递减. 因为 \(3>2>1\)，有 \(f(3)<f(2)<f(1)\). 又 \(f\) 是偶函数，故 \(f(-2)=f(2)\). 所以 \(f(3)<f(-2)<f(1)\)，故选 A.
\end{solution}

\begin{problem}
若正方体的棱长为 \(\sqrt{2}\)，则以该正方体各个面的中心为顶点的凸多面体的体积为（\quad）

\choices
    {\(\frac{\sqrt{2}}{6}\)}
    {\(\frac{\sqrt{2}}{3}\)}
    {\(\frac{\sqrt{3}}{3}\)}
    {\(\frac{2}{3}\)}
\end{problem}

\begin{answer}
B
\end{answer}

\begin{solution}
正方体中心到任一面的中心距离等于棱长的一半，即 \(r=\dfrac{\sqrt{2}}{2}\). 六个面中心构成正八面体；取通过四个侧面中心的正方形作赤道面，该正方形的边长为 \(1\)，面积为 \(1\)，正八面体可分为上下两个全等的四棱锥，每个高为 \(r\). 因而体积为 \(2\times\dfrac{1}{3}\times1\times\dfrac{\sqrt{2}}{2}=\dfrac{\sqrt{2}}{3}\)，故选 B.
\end{solution}

\begin{problem}
设曲线 \( y = x^{n + 1} \)（\(n \in \N^*\)）在点 \((1,1)\) 处的切线与 \(x\) 轴的交点的横坐标为 \( x_n \)，则 \( x_1 \cdot x_2 \cdots x_n \) 的值为（\quad）

\choices
    {\(\frac{1}{n}\)}
    {\(\frac{1}{n + 1}\)}
    {\(\frac{n}{n + 1}\)}
    {\(1\)}
\end{problem}

\begin{answer}
B
\end{answer}

\begin{solution}
曲线 \(y=x^{n+1}\) 在 \((1,1)\) 处的导数为 \(y'=(n+1)x^{n}\)，所以切线斜率为 \(n+1\)，切线方程为 \(y-1=(n+1)(x-1)\)，即 \(y=(n+1)x-n\). 令 \(y=0\)，得 \(x_{n}=\dfrac{n}{n+1}\). 因而
\[
x_{1}x_{2}\cdots x_{n}=\prod_{k=1}^{n}\dfrac{k}{k+1}=\dfrac{1}{2}\cdot\dfrac{2}{3}\cdots\dfrac{n}{n+1}=\dfrac{1}{n+1}
\]
故选 B.
\end{solution}

\section{填空题}

\begin{problem}
设等差数列 \(\{a_{n}\}\) 的前 \(n\) 项和为 \(S_{n}\)，若 \(a_{6} = S_{3} = 12\)，则 \(\{a_{n}\}\) 的通项 \(a_{n} =\fillinblank{}\).
\end{problem}

\begin{answer}
\(a_{n} = 2n\)
\end{answer}

\begin{solution}
设等差数列首项为 \(a_{1}=a\)，公差为 \(d\). 由 \(a_{6}=12\) 得 \(a+5d=12\). 由 \(S_{3}=12\) 得 \(\dfrac{3(a_{1}+a_{3})}{2}=3(a+d)=12\)，即 \(a+d=4\). 两式相减得 \(4d=8\)，所以 \(d=2\)，再得 \(a=2\). 因而 \(a_{n}=a+(n-1)d=2+2(n-1)=2n\).
\end{solution}

\begin{problem}
设 \(x\)，\(y\) 满足约束条件 \(\begin{cases}
x + y \geqslant 1,\\
x - y \geqslant -1,\\
2x - y \leqslant 2,
\end{cases}\) 目标函数 \(z = x + 2y\) 的最小值是\fillinblank{}，最大值是\fillinblank{}.
\end{problem}

\begin{answer}
\(1,11\)
\end{answer}

\begin{solution}
约束条件分别给出半平面 \(y\geqslant1-x\)、\(y\leqslant x+1\)、\(y\geqslant2x-2\). 可行域的边界交点为：前两条直线交于 \((0,1)\)，第一、第三条直线交于 \((1,0)\)，第二、第三条直线交于 \((3,4)\). 目标函数是线性函数，其最值在可行域的顶点取得. 代入三点得 \(z(0,1)=2\)、\(z(1,0)=1\)、\(z(3,4)=11\)，所以最小值为 \(1\)，最大值为 \(11\).
\end{solution}

\begin{problem}
如图球 \(O\) 的半径为 \(2\)，圆 \(O_{1}\) 是一小圆，\(O_{1}O = \sqrt{2}\). \(A\)，\(B\) 是圆 \(O_{1}\) 上两点. 若 \(A\)，\(B\) 两点间的球面距离为 \(\frac{2\pi}{3}\)，则 \(\angle AO_{1}B =\fillinblank{}\).

\bitmapfigure[max width=0.34\linewidth]{img/2009/shaanxi_liberal/q15_fig1.png}
\end{problem}

\begin{answer}
\(\frac{\pi}{2}\)
\end{answer}

\begin{solution}
设球面上两点 \(A\)，\(B\) 所对的球心角为 \(\angle AOB=\theta\). 球面距离等于半径乘球心角，所以 \(2\theta=\dfrac{2\pi}{3}\)，从而 \(\theta=\dfrac{\pi}{3}\). 在等腰三角形 \(AOB\) 中，\(OA=OB=2\)，由余弦定理得 \(AB^{2}=2^{2}+2^{2}-2\cdot2\cdot2\cos\dfrac{\pi}{3}=4\)，即 \(AB=2\). 因为小圆所在平面垂直于 \(OO_{1}\)，所以 \(OO_{1}\perp O_{1}A\)，从而 \(O_{1}A^{2}=OA^{2}-OO_{1}^{2}=4-2=2\)，同理 \(O_{1}B=\sqrt{2}\). 设 \(\angle AO_{1}B=\varphi\)，则在三角形 \(AO_{1}B\) 中由余弦定理得 \(AB^{2}=2+2-4\cos\varphi\). 代入 \(AB^{2}=4\) 得 \(\cos\varphi=0\)，且角为三角形内角，故 \(\varphi=\dfrac{\pi}{2}\).
\end{solution}

\begin{problem}
某班有 \(36\) 名同学参加数学，物理，化学课外探究小组，每名同学至多参加两个小组，已知参加数学，物理，化学小组的人数分别为 \(26\)，\(15\)，\(13\)，同时参加数学和物理小组的有 \(6\) 人，同时参加物理和化学小组的有 \(4\) 人，则同时参加数学和化学小组的有\fillinblank{}人.
\end{problem}

\begin{answer}
8
\end{answer}

\begin{solution}
设同时参加数学、化学小组的人数为 \(x\). 因为每名同学至多参加两个小组，所以三组没有公共交集. 由容斥原理，三组并集人数为
\[
26+15+13-6-4-x=44-x
\]
该并集就是参加至少一个小组的 \(36\) 名同学，因此 \(44-x=36\)，解得 \(x=8\).
\end{solution}

\section{解答题}

\begin{problem}
\begin{enumerate}
\item 已知函数 \( f(x) = A \sin (\omega x + \varphi) \)，\( x \in \R \)（其中 \(A > 0\)，\(\omega > 0\)，\(0 < \varphi < \frac{\pi}{2}\)）的周期为 \( \pi \)，且图象上一个最低点为 \( M\left(\frac{2\pi}{3}, -2\right) \). 求 \( f(x) \) 的解析式；
\item 当 \(x \in \left[0, \frac{\pi}{12}\right]\) 时，求 \(f(x)\) 的最值.
\end{enumerate}
\end{problem}

\begin{solution}
\begin{enumerate}
\item 由函数的周期为 \(\pi\)，得 \(\frac{2\pi}{\omega}=\pi\)，所以 \(\omega=2\). 又因 \(M\) 是最低点且 \(A>0\)，有 \(A=2\)，并且 \(\sin\left(\frac{4\pi}{3}+\varphi\right)=-1\). 因此 \(\frac{4\pi}{3}+\varphi=\frac{3\pi}{2}+2k\pi\). 结合 \(0<\varphi<\frac{\pi}{2}\)，得 \(k=0\)，从而 \(\varphi=\frac{\pi}{6}\)，故 \(f(x)=2\sin\left(2x+\frac{\pi}{6}\right)\).
\item 令 \(t=2x+\frac{\pi}{6}\)，当 \(x\in\left[0,\frac{\pi}{12}\right]\) 时，\(t\in\left[\frac{\pi}{6},\frac{\pi}{3}\right]\). 因为正弦函数在 \(\left[\frac{\pi}{6},\frac{\pi}{3}\right]\) 上单调递增，所以 \(f(x)\) 的最小值为 \(2\sin\frac{\pi}{6}=1\)，最大值为 \(2\sin\frac{\pi}{3}=\sqrt{3}\).
\end{enumerate}
\end{solution}

\begin{problem}
\begin{enumerate}
\item 据统计，某食品企业一个月内被消费者投诉的次数为 \(0\)，\(1\)，\(2\) 的概率分别为 \(0.4\)，\(0.5\)，\(0.1\). 求该企业在一个月内共被消费者投诉不超过 \(1\) 次的概率；
\item 假设一月份与二月份被消费者投诉的次数互不影响，求该企业在这两个月内共被消费者投诉 \(2\) 次的概率.
\end{enumerate}
\end{problem}

\begin{solution}
\begin{enumerate}
\item 设一个月内被投诉的次数为随机变量 \(X\)，则 \(P(X\leqslant 1)=P(X=0)+P(X=1)=0.4+0.5=0.9\).
\item 设一月份、二月份被投诉的次数分别为 \(X\)、\(Y\). 由两个月互不影响，\(X\)、\(Y\) 相互独立，因此 \(P(X+Y=2)=P(X=0)P(Y=2)+P(X=1)P(Y=1)+P(X=2)P(Y=0)=0.4\times0.1+0.5\times0.5+0.1\times0.4=0.33\).
\end{enumerate}
\end{solution}

\begin{problem}
\begin{enumerate}
\item 如图，在直三棱柱 \(ABC - A_{1}B_{1}C_{1}\) 中，\(AB = 1\)，\(AC = AA_{1} = \sqrt{3}\)，\(\angle ABC = 60^{\circ}\). 证明： \( AB \perp A_{1}C \)；
\item 求二面角 \(A - A_{1}C - B\) 的大小.
\end{enumerate}

\bitmapfigure[max width=0.34\linewidth]{img/2009/shaanxi_liberal/q19_fig1.png}
\end{problem}

\begin{solution}
\begin{enumerate}
\item 在 \(\triangle ABC\) 中，由余弦定理得 \(AC^{2}=AB^{2}+BC^{2}-2AB\cdot BC\cos\angle ABC\)，即 \(3=1+BC^{2}-BC\)，所以 \(BC^{2}-BC-2=0\). 因为 \(BC>0\)，故 \(BC=2\). 于是 \(BC^{2}=AB^{2}+AC^{2}=1+3\)，从而 \(AB\perp AC\). 又因为三棱柱是直三棱柱，\(AA_{1}\perp\text{平面 }ABC\)，所以 \(AB\perp AA_{1}\). 直线 \(AC\)、\(AA_{1}\) 是平面 \(A_{1}AC\) 内相交的两条直线，且都垂直于 \(AB\)，因此 \(AB\perp\text{平面 }A_{1}AC\)，从而 \(AB\perp A_{1}C\).
\item 由（1）可建立直角坐标系，取 \(A(0,0,0)\)，\(B(1,0,0)\)，\(C(0,\sqrt{3},0)\)，\(A_{1}(0,0,\sqrt{3})\). 设二面角 \(A-A_{1}C-B\) 的平面角为 \(\theta\). 平面 \(A_{1}AC\) 的一个法向量为 \(\bs{n}_{1}=(1,0,0)\)，而 \(\overrightarrow{A_{1}C}=(0,\sqrt{3},-\sqrt{3})\)，\(\overrightarrow{A_{1}B}=(1,0,-\sqrt{3})\)，故平面 \(A_{1}BC\) 的一个法向量可取 \(\bs{n}_{2}=\overrightarrow{A_{1}C}\times\overrightarrow{A_{1}B}=(-3,-\sqrt{3},-\sqrt{3})\). 因此 \(\cos\theta=\frac{|\bs{n}_{1}\cdot\bs{n}_{2}|}{|\bs{n}_{1}|\,|\bs{n}_{2}|}=\frac{3}{\sqrt{15}}=\frac{\sqrt{15}}{5}\)，所以二面角的大小为 \(\arccos\frac{\sqrt{15}}{5}\).
\end{enumerate}
\end{solution}

\begin{problem}
\begin{enumerate}
\item 已知函数 \(f(x) = x^3 - 3ax - 1\)，\(a \neq 0\). 求 \( f(x) \) 的单调区间；
\item 若 \( f(x) \) 在 \(x = -1\) 处取得极值，直线 \(y = m\) 与 \(y = f(x)\) 的图象有三个不同的交点，求 \(m\) 的取值范围.
\end{enumerate}
\end{problem}

\begin{solution}
\begin{enumerate}
\item \(f'(x)=3x^{2}-3a=3(x^{2}-a)\). 当 \(a<0\) 时，\(f'(x)>0\) 对任意 \(x\in\R\) 成立，所以 \(f(x)\) 在 \(\R\) 上单调递增；当 \(a>0\) 时，\(f'(x)>0\) 的区间为 \(\left(-\infty,-\sqrt{a}\right)\) 和 \(\left(\sqrt{a},+\infty\right)\)，\(f'(x)<0\) 的区间为 \(\left(-\sqrt{a},\sqrt{a}\right)\)，所以 \(f(x)\) 在前两个区间上单调递增，在中间区间上单调递减.
\item 因为 \(f(x)\) 在 \(x=-1\) 处取得极值，且 \(f\) 可导，所以 \(f'(-1)=0\)，即 \(3-3a=0\)，得 \(a=1\). 此时 \(f(x)=x^{3}-3x-1\)，且 \(f'(x)=3(x^{2}-1)\)，所以 \(f(x)\) 在 \(\left(-\infty,-1\right)\)、\((1,+\infty)\) 上单调递增，在 \((-1,1)\) 上单调递减. 又 \(f(-1)=1\)，\(f(1)=-3\)，因此直线 \(y=m\) 与曲线有三个不同交点当且仅当 \(-3<m<1\)，故 \(m\in(-3,1)\).
\end{enumerate}
\end{solution}

\begin{problem}
\begin{enumerate}
\item 已知数列 \(\{a_{n}\}\) 满足 \(a_1 = 1\)，\(a_2 = 2\)，\(a_{n+2} = \frac{a_n + a_{n+1}}{2}\)，\(n \in \N^*\). 令 \( b_{n} = a_{n+1} - a_{n} \)，证明：\( \{b_{n}\} \) 是等比数列；
\item 求 \(\{a_{n}\}\) 的通项公式.
\end{enumerate}
\end{problem}

\begin{solution}
\begin{enumerate}
\item 由递推关系，\(b_{n+1}=a_{n+2}-a_{n+1}=\frac{a_{n}+a_{n+1}}{2}-a_{n+1}=\frac{a_{n}-a_{n+1}}{2}=-\frac{1}{2}b_{n}\). 又 \(b_{1}=a_{2}-a_{1}=1\)，所以 \(\{b_{n}\}\) 是首项为 \(1\)、公比为 \(-\frac{1}{2}\) 的等比数列，即 \(b_{n}=\left(-\frac{1}{2}\right)^{n-1}\).
\item 由 \(a_{n}=a_{1}+\sum_{k=1}^{n-1}b_{k}\)，得 \(a_{n}=1+\sum_{k=1}^{n-1}\left(-\frac{1}{2}\right)^{k-1}=1+\frac{1-\left(-\frac{1}{2}\right)^{n-1}}{1+\frac{1}{2}}=\frac{5}{3}-\frac{2}{3}\left(-\frac{1}{2}\right)^{n-1}\)，其中 \(n\in\N^*\).
\end{enumerate}
\end{solution}

\begin{problem}
\begin{enumerate}
\item 已知双曲线 \(C\) 的方程为 \( \frac{y^{2}}{a^{2}} - \frac{x^{2}}{b^{2}} = 1 \)（\(a > 0\)，\(b > 0\)），离心率 \( e = \frac{\sqrt{5}}{2} \)，顶点到渐近线的距离为 \( \frac{2\sqrt{5}}{5} \)；
\item 如图，\(P\) 是双曲线 \(C\) 上一点，\(A\)，\(B\) 两点在双曲线 \(C\) 的两条渐近线上，且分别位于第一、二象限，若 \(\overrightarrow{AP}=\lambda\overrightarrow{PB}\)，\(\lambda\in\left[\frac{1}{3},2\right]\)，求 \(\triangle AOB\) 面积的取值范围.
\end{enumerate}

\bitmapfigure[max width=0.34\linewidth]{img/2009/shaanxi_liberal/q22_fig1.png}
\end{problem}

\begin{solution}
\begin{enumerate}
\item 设双曲线的半焦距为 \(c\)，则 \(c^{2}=a^{2}+b^{2}\). 由离心率 \(e=\frac{c}{a}=\frac{\sqrt{5}}{2}\)，得 \(1+\frac{b^{2}}{a^{2}}=\frac{5}{4}\)，所以 \(b=\frac{a}{2}\). 双曲线的渐近线为 \(y=\pm\frac{a}{b}x\). 取顶点 \(V(0,a)\)，则 \(V\) 到渐近线 \(y=\frac{a}{b}x\) 的距离为 \(\frac{ab}{\sqrt{a^{2}+b^{2}}}\). 代入 \(b=\frac{a}{2}\) 得该距离为 \(\frac{a}{\sqrt{5}}\)，故 \(\frac{a}{\sqrt{5}}=\frac{2\sqrt{5}}{5}=\frac{2}{\sqrt{5}}\)，从而 \(a=2\)，\(b=1\)，所以双曲线 \(C\) 的方程为 \(\frac{y^{2}}{4}-x^{2}=1\).
\item 双曲线的两条渐近线为 \(y=2x\) 和 \(y=-2x\). 设 \(A=(u,2u)\)，\(B=(-v,2v)\)，其中 \(u>0\)，\(v>0\). 由 \(\overrightarrow{AP}=\lambda\overrightarrow{PB}\) 得 \(\overrightarrow{OP}=\frac{\overrightarrow{OA}+\lambda\overrightarrow{OB}}{1+\lambda}\)，所以 \(P\) 的坐标为 \(\left(\frac{u-\lambda v}{1+\lambda},\frac{2u+2\lambda v}{1+\lambda}\right)\). 因为 \(P\) 在双曲线 \(C\) 上，有 \(\frac{(u+\lambda v)^{2}-(u-\lambda v)^{2}}{(1+\lambda)^{2}}=1\)，即 \(uv=\frac{(1+\lambda)^{2}}{4\lambda}\). 反之，对任意 \(\lambda\in\left[\frac{1}{3},2\right]\)，任取 \(u>0\)，并令 \(v=\frac{(1+\lambda)^{2}}{4\lambda u}>0\)，则上述 \(A\)，\(B\) 和 \(P\) 满足题设条件，所以这些 \(\lambda\) 均可取到. 因此 \(S_{\triangle AOB}=\frac{1}{2}|x_{A}y_{B}-y_{A}x_{B}|=2uv=\frac{(1+\lambda)^{2}}{2\lambda}\). 记 \(S(\lambda)=\frac{(1+\lambda)^{2}}{2\lambda}\)，则 \(S'(\lambda)=\frac{\lambda^{2}-1}{2\lambda^{2}}\)，所以 \(S\) 在 \(\left[\frac{1}{3},1\right]\) 上递减，在 \([1,2]\) 上递增. 又 \(S(1)=2\)，\(S\left(\frac{1}{3}\right)=\frac{8}{3}\)，\(S(2)=\frac{9}{4}\)，故 \(\triangle AOB\) 面积的取值范围为 \(\left[2,\frac{8}{3}\right]\).
\end{enumerate}
\end{solution}
